This chapter provides the technical background for the work presented in
the remainder of this dissertation. The material covers the general
class of defect-dynamics side-channels, the specific FPGA case and its
measurement requirements, the device physics, architecture, and
reliability infrastructure of commercial ReRAM, and finally the entropy
framework used to evaluate stochastic write timing.

Section~\ref{sec:bg-sidechannels} defines hardware side-channels and
introduces \emph{pentimenti} as a recently recognized class of leakage
channel that allows for the exploitation of BTI-driven defect dynamics
in CMOS transistors. Section~\ref{sec:bg-bti} reviews the physics of
bias temperature instability (BTI), including its reaction-diffusion and
charge-trapping/detrapping dynamics, recovery characteristics, and
propagation-delay impact. Section~\ref{sec:bg-tdc} describes the
fine-grained time-to-digital converter (TDC) sensors used to extract
pentimenti from FPGA routing logic, the role of phase-locked loop (PLL)
phase shift granularity in determining sensor resolution, and the
vendor-specific dependencies that have historically constrained FPGA
pentimenti measurement to Xilinx products.

Section~\ref{sec:bg-oxram} describes the device structure and switching
physics of oxide-based ReRAM (OxRAM), covering bipolar switching,
conductive filament formation and rupture, and the resistance states
that encode stored data. Section~\ref{sec:bg-pvloops} describes the
program-verify loop used to write OxRAM cells, explains the sources of
write latency variability, and establishes why write latency is an
observable that carries information about internal device state.
Section~\ref{sec:bg-crossbar} describes the crossbar array architecture
used in commercial OxRAM memories, covering wordline and bitline
organization, the serialization constraints imposed by bipolar
switching, and the effect of these constraints on word-level write
latency. Section~\ref{sec:bg-defects} covers defect dynamics and wearout
in OxRAM devices, describing how repeated cycling degrades switching
behavior and how wearout manifests in write latency distributions.
Section~\ref{sec:bg-ecc} provides background on linear block codes and
error-correcting codes as used in commercial ReRAM memories, covering
the generator matrix, parity submatrix, syndrome decoding, and code
strength. Section~\ref{sec:bg-entropy} covers entropy sources and true
random number generators, defining min-entropy, describing the NIST SP
800-90B evaluation framework, and introducing the statistical test
suites used to validate conditioned output.